% Generated draft: requires mathematical and source review.
\documentclass[11pt]{article}
\usepackage[margin=25mm]{geometry}
\usepackage{fontspec}
\setmainfont{Latin Modern Roman}
\usepackage{amsmath,amssymb,mathtools,unicode-math}
\setmathfont{Latin Modern Math}
\usepackage{enumitem,needspace,xurl,hyperref}
\setlength{\parindent}{0pt}
\setlength{\parskip}{5pt}
\setlength{\emergencystretch}{3em}
\newcommand{\N}{\mathbb{N}}
\newcommand{\Z}{\mathbb{Z}}
\newcommand{\Q}{\mathbb{Q}}
\newcommand{\R}{\mathbb{R}}
\newcommand{\C}{\mathbb{C}}
\newcommand{\F}{\mathbb{F}}
\newcommand{\A}{\mathbb{A}}
\newcommand{\PP}{\mathbb{P}}
\newcommand{\E}{\mathbb{E}}
\newcommand{\eps}{\varepsilon}
\newcommand{\dd}{\,\mathrm{d}}
\begin{document}
\section*{UnsolvedMath ID: 700823}
\textbf{Source problem code:} catalogue entry
\subsection{Definitions and mathematical statement}
Let \(p\) be a prime. A pro-\(p\) group is a topological group isomorphic to an inverse limit of finite \(p\)-groups. A family \(\mathcal A\) of normal subgroups is pairwise incomparable if, for distinct \(N,M\in\mathcal A\), neither \(N\subseteq M\) nor \(M\subseteq N\) holds. The supplied source calls a \(p\)-group thin if every such family has at most \(p+1\) members.

To extend this definition to pro-\(p\) groups, the meaning of normal subgroup must be specified. Under the interpretation using closed normal subgroups, thinness means
\[
\forall\mathcal A\subseteq\{N\leq G:N\trianglelefteq G\text{ and }N\text{ is closed}\},\qquad
\bigl(\mathcal A\text{ is pairwise incomparable}\bigr)\Longrightarrow |\mathcal A|\leq p+1.
\]
For each prime \(p\), let \(\mathcal T_p\) denote the collection of topological isomorphism classes of thin pro-\(p\) groups under this interpretation. A precise candidate formulation of the source question is:
\[
\text{Is }|\mathcal T_p|<\infty\text{ for every prime }p?
\]
Here the finite bound may depend on \(p\). No assumption of finiteness, infinitude, or finite generation of the groups is added.

\emph{Unresolved source qualifications.} The supplied wording does not explicitly specify closed normal subgroups, counting up to topological isomorphism, or whether \(p\) is fixed when counting. Thus the displayed question is conditional on those interpretations. Using all abstract normal subgroups instead gives a different thinness condition; counting across all primes gives a different finiteness question.
\subsection{Short English statement}
For a fixed prime \(p\), a thin group has no family of more than \(p+1\) pairwise incomparable normal subgroups. The question asks whether there are only finitely many thin pro-\(p\) groups. The proposed precise formulation counts topological isomorphism classes and uses closed normal subgroups; these conventions require confirmation from the source.
\subsection{Sources}
\begin{itemize}
\item Supplied local UnsolvedMath registry record, permanent ID \(700823\), exact source code \texttt{catalogue entry}, titled \emph{Kourovka Notebook Problem 12.3}; upstream revision \texttt{d371ee917ac9647fbb06cd6fc529deec6f3dda82}.
\item The supplied record attributes the problem to R. Brandl, proposed in 1992, and identifies the source as \emph{The Kourovka Notebook (Unsolved Problems in Group Theory)}, 12th issue (1992), original problem number 12.3, with reference to the 21st issue edition (2026), \texttt{arXiv:1401.0300v48}. The supplied collection description names E. I. Khukhro and V. D. Mazurov as editors.
\item The registry labels the problem \texttt{open} and records an update timestamp of \texttt{2026-10-16T14:00:00Z}. This is supplied registry metadata, not an independently verified mathematical or literature assessment. The cited source has not been independently consulted for this draft.
\end{itemize}
\end{document}
